\section{Conclusion}
\label{sec:conclusion}

We introduced a self-evolving framework for unified visual understanding and image generation that uses only unlabeled images and no external reward model. By decomposing a frozen backbone into three LoRA-based roles (Proposer, Solver, and Generator) and deriving all training signals from internal consistency, the framework achieves +1.9 to +3.6 on visual understanding benchmarks and +3\% GenEval across BLIP3o-8B, BAGEL-7B, and VARGPT-v1.1, with one matched training recipe and no architecture-specific tuning. Ablations confirm that both visual understanding signals (STE, self-consistency) and both image generation signals (QA fidelity, cycle-consistency) contribute independently, while loop coupling analysis shows that the Solver's dual role produces genuine cross-task synergy exceeding either loop in isolation.

Image generation supervision is inherently bounded by Solver quality, and GenEval does not fully capture perceptual dimensions such as texture fidelity or aesthetic preference. These results suggest that unsupervised self-evolution provides a principled and generalizable pathway toward improving unified models, with longer training horizons, stronger internal evaluators, and extension to video and 3D settings as natural next steps.

% =====================================================================
% COMMENTED-OUT: Extended conclusion content for journal version.
% =====================================================================
%
% \paragraph{Broader implications.}
% The self-evolving framework demonstrates that post-training improvement
% is achievable without additional human annotation, which has implications
% for democratizing model improvement: teams without access to large-scale
% annotation pipelines can still improve model capabilities. However, this
% same property means that the framework could be applied without oversight
% to models with biased pretraining data, potentially amplifying existing
% biases. Future work should investigate bias propagation dynamics under
% self-evolving training.
%
% Additional experiment ideas for future work:
% - Scaling analysis (2x, 5x, 10x steps) to show saturation curves
% - Data diversity ablation (domain-specific subsets)
% - Adapter rank ablation {4, 8, 16, 32, 64}
% - Prompt perturbation count N in {3, 5, 7, 9, 11}
% - Additional generation benchmarks (DPG-Bench, T2I-CompBench)
% - Human evaluation study comparing base vs. self-evolved generations
% - Cross-task correlation analysis (understanding category vs. GenEval subcategory)
